\scanpage{59}
\mspage{29}
preted as the \inserted{\unsure{compatible with|topological|let us topologize}} group of $U$-automorphisms of the scheme $\widetilde U$, and $E_U\longrightarrow E_K$ as the quotient homomorphism induced on $\bar K$, \unsure{identified|is identified} with the algebraic closure of $K$ in $H^0(\widetilde U,\mathcal O_{\widetilde U})$, so $E_U$ acts on $\operatorname{Spec}\bar K$ in a natural way. A section of $E_U$ over $E_K$ is therefore an action of $E_K$ on $\widetilde U$, compatible with its action on $\widetilde U$, and continuous in a suitable sense (probably the \unsure{pointwise-discrete} convergence on the \deleted{$\unsure{\widetilde U_i}$} $\bar U_i$ finite over $\bar U$ between $\widetilde U$ and $\bar U$\ldots).

Let us therefore consider the

\uline{Conjecture B} Every such action of $\Gamma$ on $\widetilde U$ admits \uline{in $X'$}$=\widehat{\widetilde U}$ a fixed point and only one.
\margin{\figure{\begin{tikzpicture}\draw (0,0)--(0,-3.1);\end{tikzpicture}}}
\margin{and $\unsure{=|\overline m}$ the induced action of $\Gamma^\circ$ must have a fixed point and only one --- see below.}

This therefore means

a) If there is a fixed point at finite distance that is $\widetilde x\in\widetilde U^\Gamma$, then 1°) The image of $\widetilde x$ in $U$ is uniquely determined --- this is essentially theorem 3 of $\S1$ (two distinct points of $U(K)$ define distinct conjugacy classes of splittings) and
\margin{58}
